\documentclass[11pt,a4paper]{article}
\usepackage[utf8]{inputenc}
\usepackage[T1]{fontenc}
\usepackage[portuguese]{babel}
\usepackage{xcolor}
\usepackage{hyperref}
\usepackage{geometry}
\usepackage{tcolorbox}

\geometry{margin=2.5cm}

\definecolor{primary}{RGB}{41, 128, 185}
\definecolor{secondary}{RGB}{44, 62, 80}
\definecolor{codebg}{RGB}{240, 240, 240}

\hypersetup{colorlinks=true, linkcolor=primary, urlcolor=primary}

\title{\color{secondary}\Huge\textbf{Exercício 5.5: Comparação de Plataformas}}
\author{\color{primary}DevOps com Kubernetes -- 2026}
\date{}

\begin{document}
\maketitle

\section*{\color{primary}Instruções do Exercício}
\begin{tcolorbox}[colback=blue!5!white,colframe=blue!75!black,title=Exercício 5.5: Comparação de Plataformas]
Analise o ecossistema e as opções atuais de plataformas baseadas em Kubernetes.

Escolha um provedor de serviços/gerenciador de cluster (como o Rancher, por exemplo) e compare-o diretamente com outro (como o OpenShift ou Anthos GKE). 

Decida de maneira arbitrária qual provedor de serviço é 'melhor' para uma equipe fictícia e apresente argumentos em defesa de sua escolha contra a do outro provedor (destaque pontos fortes, modelos de uso, licença ou governança).

Para a submissão deste exercício, uma lista em tópicos (\textit{bullet points}) bem elaborada é o suficiente.
\end{tcolorbox}

\end{document}
